\documentclass[11pt,a4paper]{article}

\usepackage{amsmath}
\usepackage{amssymb}
\usepackage{amsthm}
\usepackage{graphicx}
\usepackage{hyperref}
\usepackage{booktabs}
\usepackage{algorithm}
\usepackage{algpseudocode}
\usepackage{listings}
\usepackage{xcolor}
\usepackage{subfiles}

% Theorem-like environments
\newtheorem{theorem}{Theorem}
\newtheorem{lemma}[theorem]{Lemma}
\newtheorem{proposition}[theorem]{Proposition}
\newtheorem{corollary}[theorem]{Corollary}
\newtheorem{definition}{Definition}
\newtheorem{remark}[theorem]{Remark}

\newcommand{\Ber}[1]{\widehat{#1}}
% two parameter version of Ber
\newcommand{\Berr}[2]{\widehat{#1}^{(#2)}}
\newcommand{\Bool}{\mathrm{Bool}}
\newcommand{\True}{\mathrm{True}}
\newcommand{\False}{\mathrm{False}}


\title{The Bernoulli Model: A Probabilistic Framework for Data Structures and Types}
\author{Alexander Towell\\\texttt{atowell@siue.edu}}
\date{}

\begin{document}

\maketitle

\begin{abstract}
This paper introduces the Bernoulli Model, a probabilistic framework designed to handle uncertainty in data types, particularly Boolean values. The Bernoulli Model provides a formalism for optimizing space and accuracy in data representation through controlled trade-offs. We explore its applications in developing space-efficient data structures, such as Bloom filters and Count-Min sketches, and introduce the concept of Cipher Data Types. The framework allows for $\mathcal{O}(1)$ time complexity while maintaining probabilistic guarantees on accuracy. This work forms a foundational part of a larger project aimed at enhancing data structure efficiency and error propagation analysis.
\end{abstract}

\tableofcontents

\subfile{sections/intro}
\subfile{sections/order_and_channels}
\subfile{sections/prediction}
\subfile{sections/lifting}
\subfile{sections/set_indicator}
\subfile{sections/logic_gates}
\subfile{sections/bernoulli_maps}

\bibliographystyle{plain}
\bibliography{references}

\end{document}
